\documentclass[a4paper,12pt]{article}
\usepackage{fancyhdr}
\pagestyle{fancy}
\usepackage{graphicx}
\usepackage{setspace}
\usepackage{microtype}
\usepackage{multicol}
\usepackage{multirow}
\usepackage[portuguese]{babel}


\begin{document}
\setlength{\headsep}{70pt}
\rhead{\large{Produção de Documentos Técnicos \\ Exercícios Adicionais - Modo Texto}}
\lhead{\includegraphics[width=5cm]{fculogo}} 

\begin{description}
\item A - Reproduza o texto tal como lhe é apresentado. Não se esqueça de ter atenção a todos os estilos de texto, formatações, espaçamentos, alinhamentos. Qualquer tipo de lista utilizado, deve ser feita automaticamente pelo LATEX. No final insira o Índice de Conteúdo, de Tabelas e Figuras.
\item B - 16:55
\end{description}

\section{A Descoberta dos seres Peixes}

\hspace{0.5 cm}\textit{Seres primitivos, dotados de armadura, de placas espinhosas e escamas sobrepostas, sem mandíbula e com três olhos na cabeça} deram origem, há \Large{400 milhões} \normalsize{de anos, às 30.000 espécies de peixes que conhecemos hoje.}

\singlespacing

\textbf{\Large{C\normalsize{OMO CONTAR A SUA IDADE?}}}

\singlespacing

\hspace{0.5 cm} As escamas apresentam-se com os aspetos mais variados, mas se examinarmos qualquer delas com atenção, veremos uma série de círculos; através deles podemos saber a idade dos peixes. Mas, a cada círculo não corresponde um ano de crescimento.

\hspace{0.5 cm} \textbf{Um peixe não cresce sempre o mesmo através do ano}. Nos \texttt{meses quentes}, quando a comida não falta, cresce rapidamente, o que provoca em cada escama uma quantidade de anéis bastante distanciados. Nos \texttt{meses frios}, os peixes crescem muito devagar e às vezes nem chegam a crescer. Os seus anéis de crescimento são poucos e muito fundos. Este conjunto de anéis forma uma espécie de linha mais acentuada. Para conhecermos com exatidão a idade de um peixe \underline{devemos contar apenas linhas} pois, uma, marca a passagem de um \scalebox{2}{\textls[5]{inverno}}.

\begin{itemize}
\item Exemplos de Peixes:
\begin{enumerate}
\item Carapau
\item Cavala
\item Cherne
\item Congro
\end{enumerate}
\end{itemize}
\begin{itemize}
\item Exemplos de Cefalópodes:
\begin{description}
\item[--] Choco
\item[--] Lula
\item[--] Polvo-comum
\end{description}
\end{itemize}

\newpage
\pagestyle{plain}

\begin{table}[t]
\caption{Comparação entre Peixes Cartilaginosos e Ósseos}
\begin{tabular}{ c || cc }
\multicolumn{2}{c}{} \\ \hline
\textbf{CARACTERÍSTICAS} & \textbf{CARTILAGINOSOS} & \textbf{ÓSSEOS} \\ \hline \hline
\texttt{Esqueleto} & Cartilaginoso & Ósseo \\ 
\texttt{Escamas} & Placóides & Ausentes ou Ciclóides \\
\texttt{Válvula Espiral} & Ausente & Presente \\
\texttt{Brânquias} & Sem Opérculo & Com Opérculo  \\ \hline
\end{tabular}
\end{table}

\section{Valor Nutricional do Peixe}

\begin{table}[ht]
\begin{tabular}{| cc | c |}
\multicolumn{2}{c}{} \\ \hline
\multirow{5}{*} 
Peixes  &               & Qtd. a ser ingerida \\ \hline
{Atum} & enlatado & 340 gr \\ \hline
            & fresco & 71-340 gr \\ \hline
{Salmão} & rosa & 71 gr \\ \hline
               & do Atlântico (cultivado) & 42.5-71 gr \\ \hline
               & do Atlântico (selvagem) & 57-99 gr \\ \hline
\end{tabular}
\caption{Valor Nutricional do Peixe}
\end{table} 

\begin{center}
\begin{figure}[h!]
\scalebox{0.6}{\includegraphics{Robalo}}
\caption{Um Robalo}
\end{figure}
\end{center}

\newpage

\tableofcontents
\listoffigures
\listoftables


\end{document}